\section{药物跨膜转运}

\begin{frame}
    \frametitle{药物跨膜转运}
    药物在体内吸收、分布、转化及排泄的过程中，首先必须跨越多层生物膜，进行多次转运，这种过程叫做药物跨膜转运。因此，了解药物跨膜转运机制及其影响因素是十分重要的。
\end{frame}

\subsection{生物膜的结构}

\begin{frame}
    \frametitle{经典模型}
    
    \begin{quote}
        生物膜=磷脂质+蛋白质+糖类
    \end{quote}

    1935年，Danielli与Davson提出细胞膜经典模型。
    
\begin{itemize}
    \item 细胞膜是由脂质双分子层构成；
    \item 两个脂质分子的疏水尾相连在中间形成膜的疏水区；
    \item 脂质分子的亲水头朝外分布在膜的外侧形成对称的双层膜结构；
    \item 膜蛋白分布在脂质双分子层的两侧；
    \item 膜上有许多带电荷的小孔，水分能自由通过；
    \item 膜结构中还存在特殊载体和酶促系统，能与某些物质特异性结合，进行物质转运。
  \end{itemize}
\end{frame}

\begin{frame}
    \frametitle{流动镶嵌模型}
    \hspace{-2em}
    \begin{minipage}{0.38\linewidth}

        1972年，Singer和Nicolson提出生物膜流动讓嵌模型，该模型强调了膜的流动性和不对称性，即膜的结构不是静止的而是流动的。

\begin{itemize}
  \item 该模型的基本结构仍是脂质双分子层
  \item 流动的脂质双分子层构成细胞膜的连续主体
  \item 蛋白质分子以不同的方式和不同的深度嵌入磷脂双分子层中
\end{itemize}

    \end{minipage}\hfill
    \begin{minipage}{0.6\linewidth}
        \begin{figure}[htpb]
            \centering
            \includegraphics[width=\linewidth]{流动镶嵌模型}
        \end{figure}
    \end{minipage}
    
\end{frame}

\begin{frame}
    \frametitle{晶格镶嵌模型}
    1975年由Wallach提出晶格镶嵌模型，进一步解释了膜的流动性和完整性特征
    \begin{itemize}
      \item 膜的流动性是由于脂质能可逆地进行无序（液态）和有序（晶态）的相变过程
      \item 膜蛋白对脂质分子的活动具有控制作用
      \item 具有流动性的脂质是呈小片的点状分布
      \item 脂质的流动性是局部的，并不是整个脂质双分子层都在流动
    \end{itemize}
    解释了为什么细胞膜既具有流动性又能保持其完整性和稳定性。
\end{frame}


\begin{frame}
    \frametitle{生物膜的性质}
    \begin{itemize}
      \item 流动性
      \item 不对称性
      
      细胞膜内外两侧层面的组分和功能有明显的差异
      \item 选择透过性
      
      细胞膜具有选择透过性 ，可以让水分子自由透过 ，选择吸收的离子和小分子也可以通过 ，而其他的离子 、小分子和大分子则不能通过 。例如 ，膜的脂质结构特征 ，使得脂溶性药物较易透过而脂溶性很小的药物难以透过 ；膜上特异性表达的载体蛋白是某些药物选择性透过的载体 。
    \end{itemize}
\end{frame}


\subsection{药物的跨膜转运方式}

\begin{frame}
    \frametitle{药物的跨膜转运方式}
\begin{figure}[htpb]
    \centering
    \includegraphics[width=.78\linewidth]{自由扩散}
\end{figure}
\end{frame}

\begin{frame}
    \frametitle{被动扩散(单纯扩散)}
    被动转运是指不需要消耗能量，生物膜两侧的药物由高浓度侧向低浓度侧转运的过程。大多数药物是通过这种方式转运的，即药物是借助于在生物膜中的脂溶性顺浓度差跨膜转运的。
    
        \begin{enumerate}
          \item 顺浓度梯度转运，可以用Fick定律描述药物的转运速率
          \[
            \frac{\mathrm{d}C}{\mathrm{d}t}=\frac{DAk}{h}(C_{\mathrm{GI}}-C)
          \]
          $\frac{\mathrm{d}C}{\mathrm{d}t}$为扩散速率；$D$ 为扩散系数；$A$为扩散表面积， $k$为油水分配系数；$h$为膜厚度 ；$C_{\mathrm{GI}}$为胃肠道中的药物浓度 ；$C$为血药浓度。
          \item 膜对转运的药物无选择性，转运速率只与其脂溶性大小及扩散系数有关
          \item 膜对转运的药物量无饱和现象
          \item 无竞争性抑制作用
        \end{enumerate}
\end{frame}


\begin{frame}
    \frametitle{被动扩散(单纯扩散)}
    \[
        \frac{\mathrm{d}C}{\mathrm{d}t}=\frac{DAk}{h}(C_{\mathrm{GI}}-C)
    \]
    当药物口服后 ，胃肠道中的浓度远大于血中的药物浓度，则$C$可以忽略不计；在给予某一药物于某一个体的吸收过程中，其 $D$、$A$、$h$、$k$都为定值 ，可用渗透系数$P$ (permeability coefficient) 来表达 ， 即$ P =DAk/h$ ， 此时
    \[
        \frac{\mathrm{d}C}{\mathrm{d}t}=PC_{\mathrm{GI}}
    \]
    即药物的扩散速率等于渗透系数与胃肠道药物浓度的乘积 。渗透系数是药物渗透性的参数，描述药物的膜渗透能力。对于给定药物，药物单纯扩散透过膜的转运速率与胃肠道中药物浓度呈线性关系 。
\end{frame}

\begin{frame}
  \frametitle{促进扩散}
  又称易化扩散、协助扩散，或帮助扩散。

  促进扩散是指某些药物在细胞膜上转运体的帮助下 ， 由高浓度侧向低浓度侧跨细胞膜转运的过程 。有些药物虽然水溶性不好 、 脂溶性也较差 ，但也能有较好的透膜吸收， 这是因为膜结构中的一些特殊转运体参与了药物转运 。 一般认为 ， 促进扩散的转运机制是 ： 细胞膜上的转运体在膜外侧与药物结合后 ， 通过转运体的自动旋转或变构将药物转运到细胞膜内侧 。

  \begin{enumerate}
    \item 速率快 、 效率高
    \item 有选择性
    \item 有饱和现象
    \item 有部位特异性
    \item 有竞争性抑制现象
  \end{enumerate}

\end{frame}


\begin{frame}
    \frametitle{孔道转运}
    生物膜上有水通道或蛋白质分子孔道。在毛细血管壁上也存在细胞间隙构成的孔道。一些物质包括水和某些电解质等可以通过这些孔道转运。这种转运通常与药物的分子结构和大小有关。转运速率主要取决于相应组织的血流速率以及生物膜的性质，而与脂溶性和pH的梯度关系不大。但脑等特殊组织，由于毛细管内皮细胞紧密联接，缺乏孔道转运。 

    \begin{enumerate}
      \item 主要为水和电解质的转运
      \item 转运速率与所处组织及膜的性质有关
    \end{enumerate}
\end{frame}


\begin{frame}
    \frametitle{主动转运}
    \begin{figure}[htpb]
      \centering
      \includegraphics[width=.9\linewidth]{主动运输}
    \end{figure}
\end{frame}


\begin{frame}
    \frametitle{主动转运}
    主动转运是指生物膜两侧的药物借助载体蛋白的帮助由低浓度侧向高浓度侧（ 逆浓度梯度 ） 转运的过程 。主动转运需要消耗能量。

    主动转运速率可用米氏方程描述：
    \[
        \frac{\mathrm{d}C}{\mathrm{d}t}=\frac{V_mC}{K_m+C}
    \]
    主动转运可分为 ATP 驱动泵和协同转运两种 。
    \begin{enumerate}
      \item ATP 驱动泵
      
      以 ATP 水解释放的能量为直接能源进行主动转运的载体蛋白家族称为ATP 驱动泵 。 这类载体蛋白也是一种 ATP 酶 ，能催化 ATP 水解提供能量 。
      \item 协同转运
      
      转运器不直接利用 ATP 水解的能量 ， 而是借助膜上相邻钠钾泵排出Na$^+$ 所产生的势能贮备，与 Na$^+$ 相伴或后继进行药物转运的主动转运方式。
    \end{enumerate}
\end{frame}


\begin{frame}
    \frametitle{主动转运特点}
    \begin{enumerate}
      \item 逆浓度梯度转运
      \item 常需要能量
      \item 需要载体参与
      \item 转运速率有饱和性
      \item 有竞争性抑制现象，同类物质往往能竞争同一载体, 产生相互拮抗作用。
      \item 对物质转运存在化学和立体选择性
    \end{enumerate}
\end{frame}


\begin{frame}
    \frametitle{膜动转运}
    转运蛋白虽然能够帮助许多离子和小的分子通过细胞膜，但是，对于像蛋白质和多糖这样的生物大分子的运输却无能为力。
    \begin{figure}[htpb]
      \centering
      \includegraphics[width=.6\linewidth]{胞吞与胞吐}
    \end{figure}
\end{frame}


\begin{frame}
    \frametitle{膜动转运}
    通过细胞膜的主动变形将物质摄入细胞内或从细胞内释放到细胞外的转运过程。
    \begin{enumerate}
      \item 入胞作用
      
      物质借助与细胞膜上某些蛋白质的特殊亲和力而附着于细胞膜上，通过细胞膜的内陷形成小泡，包裹药物的小泡逐渐与细胞膜表面断离，从而将物质摄入细胞内的转运过程

      \begin{itemize}
        \item 蛋白质 、 多肽类 、 脂溶性维生素
        \item 重金属
      \end{itemize}
      \item 出胞作用
      
      某些大分子物质通过形成小泡从细胞内部移至细胞膜内表面，小泡的膜与细胞膜融合 ，从而将物质排出细胞外的转运过程

      \begin{itemize}
        \item 腺细胞分泌胰岛素
        \item 细胞内不能消化的物质
      \end{itemize}
    \end{enumerate}
\end{frame}

\begin{frame}
    \frametitle{药物的转运机制是一个非常复杂的过程}
    药物以何种机制转运吸收 ， 与药物性质 、 吸收部位生理特征等密切相关 。 
    
    某种药物可以通过一种特定的转运机制吸收 ， 也可以通过多种转运机制吸收 。
\end{frame}


\begin{frame}[t]
    \frametitle{}
    \begin{question}
        药物借助载体或酶促系统，消耗机体能量，从膜的低浓度向高浓度一侧转运的药物转运方式是
        \tcblower

        \only<1>{
            A. 主动转运
        }
        \only<2>{
            \highlight{A. 主动转运}
        }

        B. 简单扩散

		    C. 易化扩散

        D. 膜动转运

        
    \end{question}
\end{frame}


\begin{frame}[t]
    \frametitle{}
    \begin{question}
        扩散速度取决于膜两侧药物的浓度梯度，药物的脂水分配系数及药物在膜内扩散速度的药物转运方式是
        \tcblower

        A. 滤过

        \only<1>{
            B. 简单扩散
        }
        \only<2>{
            \highlight{B. 简单扩散}
        }

        C. 易化扩散
        
        D. 主动转运

        
    \end{question}
\end{frame}